\documentclass[11pt]{amsart}
\usepackage[T1]{fontenc}
\usepackage{lmodern}
\usepackage{amsmath,amssymb,amsthm,mathtools}
\usepackage{booktabs}
\usepackage{array}
\usepackage{microtype}
\usepackage{xcolor}
\usepackage[colorlinks=true,linkcolor=blue!45!black,citecolor=blue!45!black,urlcolor=blue!45!black]{hyperref}

\newtheorem{theorem}{Theorem}[section]
\newtheorem{lemma}[theorem]{Lemma}
\newtheorem{proposition}[theorem]{Proposition}
\newtheorem{corollary}[theorem]{Corollary}
\theoremstyle{definition}
\newtheorem{definition}[theorem]{Definition}
\newtheorem{remark}[theorem]{Remark}

\newcommand{\mdim}{\operatorname{md}}
\newcommand{\diam}{\operatorname{diam}}

\title[Starlike trees and multiset dimension]{Starlike Trees and the Diameter Conjecture\\for Multiset Dimension}
\author{Zijian Zeng}
\address{Institute of Computer Science and Digital Innovation\\
UCSI University, Kuala Lumpur 56000, Malaysia}
\email{zijianzeng@foxmail.com}
\email{1002266693@ucsiuniversity.edu.my}
\date{August 30, 2026}
\subjclass[2020]{05C05, 05C12, 05C35}
\keywords{multiset dimension, multiset resolving set, ID-coloring, tree, starlike tree, diameter}
\hypersetup{
  pdftitle={Starlike Trees and the Diameter Conjecture for Multiset Dimension},
  pdfauthor={Zijian Zeng}
}

\begin{document}

\begin{abstract}
For a vertex set \(W\) in a graph, the multiset representation of \(v\) is the
multiset of its distances to the vertices of \(W\).  A conjecture of Hafidh
et al. states that every \(n\)-vertex tree of finite multiset dimension
satisfies \(\mdim(T)\le n-\diam(T)+1\).  We prove the conjecture for every
starlike tree.  More precisely, if the ordered arm lengths are
\(l_1\le\cdots\le l_m\), we show that the tree has finite multiset dimension
if and only if at most \(2^r\) arms have length at most \(r\), for every
\(r\ge1\).  The sufficiency proof assigns to the arms canonical binary
landmark patterns.  An odd-modulus distance-sum argument recovers vertex
depth, and a shell-count polynomial recovers the arm.  A binary-weight
estimate then places this resolving set within the conjectured diameter
budget.  Independently, two exact implementations verify the conjecture for
all 13,188 nonisomorphic trees of order at most 15.  The conjecture for
arbitrary trees remains open.
\end{abstract}

\maketitle

\section{Introduction}

Let \(G\) be a finite connected graph and \(W\subseteq V(G)\).  The multiset
representation of \(v\in V(G)\) with respect to \(W\) is
\[
 r_m(v\mid W)=\{\!\{d(v,w):w\in W\}\!\}.
\]
If these multisets are pairwise distinct, then \(W\) is a multiset resolving
set.  The minimum cardinality of such a set is the multiset dimension
\(\mdim(G)\); if no such set exists, \(\mdim(G)=\infty\).  This parameter was
introduced by Simanjuntak, Siagian, and Vetr\'ik~\cite{Simanjuntak2017}.
It is equivalent to the ID-number defined through distance-shell counts
\cite{HakanenYero2024}.

Hafidh et al.~\cite{Hafidh2019} proved that an \(n\)-vertex tree \(T\) of
diameter at least two and finite multiset dimension satisfies
\(\mdim(T)\le n-2\).  They proposed the sharper conjecture
\begin{equation}
 \mdim(T)\le n-\diam(T)+1,                                  \tag{C}
\end{equation}
and exhibited an infinite equality family.  The conjecture is still listed
as open in a 2026 survey~\cite{Albejani2026}.  The general-graph extremal
problem has moved in a different direction: finite multiset dimension can
equal the order of a graph~\cite{Allikvere2026}.

A \emph{starlike tree} has exactly one vertex of degree greater than two; its
components after deleting that vertex are paths.  Kono and Zhang
\cite[Proposition~2.1]{KonoZhang2021} proved a necessary capacity condition
for such a tree to admit an ID-coloring, and supplied several sufficient
families.  Our main result proves that their capacity condition is also
sufficient and, at the same time, proves (C) for the entire class.

\begin{theorem}\label{thm:intro}
Let \(T\) be a starlike tree with \(m\ge3\) arms of ordered lengths
\(1\le l_1\le\cdots\le l_m\).  Then the following are equivalent.
\begin{enumerate}
\item \(\mdim(T)<\infty\).
\item For every \(r\ge1\), at most \(2^r\) arms have length at most \(r\).
\end{enumerate}
Whenever these conditions hold,
\[
 \mdim(T)\le |V(T)|-\diam(T)+1.
\]
\end{theorem}

The proof is constructive and uses neither probabilistic choices nor
asymptotic estimates.  We also report an exact finite verification through
order 15.  These results are partial progress on (C), not a solution for
trees with several branching vertices.

\section{Shell counts and short-arm capacity}

For \(W\subseteq V(G)\), let
\[
 c_j(v;W)=|\{w\in W:d(v,w)=j\}|.
\]
The multiset \(r_m(v\mid W)\) and the sequence
\((c_0(v;W),c_1(v;W),\ldots)\) contain exactly the same information.  The
distance-zero coordinate can also be omitted without changing whether
vertices are resolved, since
\[
 c_0(v;W)=|W|-\sum_{j\ge1}c_j(v;W).
\]
This gives directly the convention-level equivalence with ID-colorings.

Let the center of a starlike tree be \(o\).  On arm \(i\), write \(v_{i,s}\)
for the vertex at depth \(s\), where \(1\le s\le l_i\).

\begin{lemma}[short-arm capacity]\label{lem:capacity}
If a starlike tree has a multiset resolving set, then for every \(r\ge1\) it
has at most \(2^r\) arms of length at most \(r\).  Consequently,
\[
 l_i\ge\lceil\log_2 i\rceil\qquad(2\le i\le m).
\]
\end{lemma}

\begin{proof}
Fix a resolving set \(W\).  Encode each arm of length at most \(r\) by the
binary word \((b_1,\ldots,b_r)\), where \(b_s=1\) exactly when its depth-\(s\)
vertex belongs to \(W\), and put \(b_s=0\) beyond the end of the arm.

Suppose that two arms have the same word.  For their depth-one vertices, the
distances to selected vertices on these two arms are exchanged: a selected
depth-\(s\) vertex on the same arm is at distance \(s-1\), and its counterpart
on the other arm is at distance \(s+1\).  Distances to the center and to every
selected vertex on all remaining arms are identical.  Hence the two
depth-one vertices have the same multiset representation, a contradiction.
Thus the short arms have distinct words, of which there are \(2^r\).
The ordered-length inequality follows by taking \(r=l_i\).
\end{proof}

\section{Canonical binary landmark patterns}

For a nonnegative integer \(q\), define
\[
 S_q=\{a\ge1:\text{the bit of weight }2^{a-1}\text{ in }q\text{ is }1\}.
\]
Thus \(S_0=\varnothing,S_1=\{1\},S_2=\{2\},S_3=\{1,2\}\), and so on.
Put
\[
 k_m=\sum_{q=0}^{m-1}|S_q|,
 \qquad
 K_m=k_m+\mathbf 1_{\{k_m\text{ is even}\}}.
\]
In particular, \(K_m\) is odd.

Assume the capacity inequalities of Lemma~\ref{lem:capacity}.  On arm \(i\),
select the vertices at the depths in \(S_{i-1}\).  This is feasible because
\(\max S_{i-1}=\lceil\log_2i\rceil\) when \(i\ge2\).  If \(k_m\) is even,
also select the center \(o\).  Denote the resulting \(K_m\)-vertex set by
\(W_m\).

\begin{lemma}\label{lem:binary-resolves}
The set \(W_m\) resolves \(T\).
\end{lemma}

\begin{proof}
Put
\[
 A=\sum_{q=0}^{m-1}\sum_{a\in S_q}a,
 \qquad
 G_q(s)=\sum_{a\in S_q}\min\{s,a\}.
\]
For a vertex \(v_{q+1,s}\) at depth \(s\ge1\) on the arm carrying \(S_q\),
the sum of the entries of its multiset representation is
\begin{equation}
 F(q,s):=\sum_{w\in W_m}d(v_{q+1,s},w)
       =K_ms+A-2G_q(s).                                    \label{eq:sum}
\end{equation}
Indeed, treating every selected arm vertex as if it were on another arm
gives \(K_ms+A\); for a selected vertex at depth \(a\) on the same arm, the
distance changes from \(s+a\) to \(|s-a|=s+a-2\min\{s,a\}\).

We need the strict estimate
\begin{equation}
 0\le G_q(s)\le\sum_{a\in S_q}a\le q\le m-1<K_m.            \label{eq:G}
\end{equation}
The middle inequality uses \(a\le2^{a-1}\).  Moreover,
\(k_m\ge m-1\), since each of \(1,\ldots,m-1\) has a set bit.  When \(m=3\),
\(k_m=2\) and the center is added.  When \(m\ge4\), the number \(3<m\) has
two set bits, so \(k_m>m-1\).  This proves the last inequality in
\eqref{eq:G}.

If two noncentral vertices have equal representations, then their distance
sums are equal.  By \eqref{eq:sum},
\[
 K_m(s-t)=2\bigl(G_q(s)-G_{q'}(t)\bigr).
\]
Because \(K_m\) is odd, it divides the difference in parentheses.  Its
absolute value is less than \(K_m\) by \eqref{eq:G}; hence the difference is
zero and \(s=t\).

It remains to distinguish arms at a common depth \(s\).  Encode a multiset of
distances by its counting polynomial.  Define
\[
 P(z)=\sum_{q=0}^{m-1}\sum_{a\in S_q}z^a.
\]
Apart from the common term \(z^s\) when the center is selected, the counting
polynomial for \(v_{q+1,s}\) is
\begin{equation}
 z^sP(z)+\sum_{a\in S_q}\bigl(z^{|s-a|}-z^{s+a}\bigr).
                                                               \label{eq:poly}
\end{equation}
For fixed \(s\), the polynomials
\(z^{|s-a|}-z^{s+a}\), \(a\ge1\), are linearly independent.  In a nonzero
finite linear combination, choose the largest \(a\) with nonzero coefficient.
The term of degree \(s+a\) can be produced neither by a smaller negative term
nor by any positive term \(z^{|s-b|}\).  Thus \eqref{eq:poly} determines
\(S_q\), and the sets \(S_q\) are pairwise distinct.

Finally, the distance sum at the center is \(A\).  Equality with
\eqref{eq:sum} would give \(K_ms=2G_q(s)\).  Oddness and \eqref{eq:G} would
force \(G_q(s)=0\), and then \(s=0\), a contradiction.  The center is also
resolved.
\end{proof}

Lemma~\ref{lem:capacity} and Lemma~\ref{lem:binary-resolves} already prove the
finiteness equivalence in Theorem~\ref{thm:intro}.  The next section controls
the number of selected vertices.

\section{The diameter budget}

The two longest arms form a diametral path, so
\begin{equation}
 |V(T)|-\diam(T)+1=2+\sum_{i=1}^{m-2}l_i.                  \label{eq:budget}
\end{equation}

\begin{lemma}[binary-weight budget]\label{lem:weight}
For \(m=3\) or \(m\ge5\),
\[
 K_m\le 2+\sum_{i=1}^{m-2}\max\{1,\lceil\log_2i\rceil\}.
\]
\end{lemma}

\begin{proof}
The assertion for \(m=3\) is \(3\le3\).  Let \(m\ge5\), and put
\[
 b(m)=2+\sum_{i=1}^{m-2}\max\{1,\lceil\log_2i\rceil\}.
\]
For \(L\ge2\), grouping the integers \(i\le L\) by binary length gives, with
\(h=\lceil\log_2L\rceil\),
\begin{equation}
 \sum_{i=1}^{L}\max\{1,\lceil\log_2i\rceil\}
 =2+hL-2^h.                                                \label{eq:blocksum}
\end{equation}

Let \(H=\lceil\log_2m\rceil\).  In the initial binary words
\(0,\ldots,m-1\), each bit is 1 no more often than it is 0: its periodic
pattern consists of an equally long zero block followed by a one block, and
every incomplete period starts with zeroes.  Hence
\begin{equation}
 k_m\le \frac{mH}{2}.                                      \label{eq:popbound}
\end{equation}

First suppose \(m\ge2^{H-1}+3\), so that
\(\lceil\log_2(m-2)\rceil=H\).  By
\eqref{eq:blocksum}--\eqref{eq:popbound}, it is enough that
\[
 \frac{mH}{2}+1\le4+H(m-2)-2^H.
\]
For \(H\ge4\), the difference is increasing in \(m\), and at
\(m=2^{H-1}+3\) it is
\[
 2^{H-2}(H-4)-\frac H2+3\ge0.
\]

It remains to treat the first two integers in a binary block.  If
\(m=2^{H-1}+1\), then
\(k_m=(H-1)2^{H-2}+1\), and \(b(m)-(k_m+1)\) equals
\((H-3)(2^{H-2}-1)\).  If \(m=2^{H-1}+2\), then
\(k_m=(H-1)2^{H-2}+3\), and the difference is
\((H-3)2^{H-2}\).  The remaining block \(H=3\), namely
\(m=5,6,7,8\), gives
\[
 (k_m+1,b(m))=(6,6),(8,8),(10,11),(13,14).
\]
Since \(K_m\le k_m+1\), the lemma follows.
\end{proof}

The exceptional value \(m=4\) needs a smaller construction than the parity
adjustment in \(W_m\).

\begin{lemma}\label{lem:four}
If \(m=4\) and the capacity condition holds, then \(T\) has a resolving set of
size at most \(l_1+l_2+2\).
\end{lemma}

\begin{proof}
Name the arms \(A,B,C,D\) in nondecreasing order of length.  The capacity
condition gives \(l_3\ge2\).  If \(l_4\ge3\), take
\[
 W=\{A_1,C_2,D_3\}.
\]
Temporarily denote the arms containing landmarks at depths \(1,2,3\) by
\(X,Y,Z\), and the remaining arm by \(Q\).  For the sum \(F\) of distances to
\(W\),
\[
\begin{array}{c|c@{\qquad}c|c}
x&F(x)&x&F(x)\\ \hline
o&6&X_i&3i+4\\
Y_1&7&Y_i\ (i\ge2)&3i+2\\
Z_1&7&Z_2&8\\
Z_i\ (i\ge3)&3i&Q_i&3i+6.
\end{array}
\]
The only equal-sum groups are \(X_1,Y_1,Z_1\), the pair \(Y_2,Z_2\), and the
pairs \(Z_{i+2},Q_i\).  Their representations are, respectively,
\[
 \{0,3,4\},\ \{1,2,4\},\ \{2,2,3\};
\]
\[
 \{0,3,5\},\ \{1,3,4\};
\]
and
\[
 \{i-1,i+3,i+4\},\ \{i+1,i+2,i+3\}.
\]
Thus \(W\) resolves.

If \(l_4\le2\), the ordered tuple is one of
\((1,1,2,2),(1,2,2,2),(2,2,2,2)\).  Resolving sets are
\[
 \{A_1,C_1,C_2,D_2\},\qquad
 \{A_1,B_1,B_2,C_2\},\qquad
 \{A_1,A_2,B_1,C_2\},                                    \tag{8}
\]
respectively.  A direct listing of the representations in the same order is
\[
\begin{array}{c|l}
(1,1,2,2)&
1122;\ 0233;\ 2233;\ 0123,0134;\ 1223,0334\\
(1,2,2,2)&
1122;\ 0233;\ 0123,0134;\ 1223,0334;\ 2233,3344\\
(2,2,2,2)&
1122;\ 0123,0134;\ 0233,1344;\ 1223,0334;\ 2233,3344.
\end{array}
\]
Here, for example, \(0233\) denotes the multiset
\(\{\!\{0,2,3,3\}\!\}\).  Every row contains pairwise distinct entries.
The sets in (8) have four vertices, while \(l_1+l_2+2\ge4\).
\end{proof}

\begin{proof}[Proof of Theorem~\ref{thm:intro}]
The equivalence follows from Lemmas~\ref{lem:capacity} and
\ref{lem:binary-resolves}.  If \(m=3\), Lemma~\ref{lem:weight} and
\eqref{eq:budget} bound the three-vertex canonical set.  If \(m\ge5\), combine
Lemma~\ref{lem:capacity}, Lemma~\ref{lem:weight}, and \eqref{eq:budget}.  The
case \(m=4\) is Lemma~\ref{lem:four}.
\end{proof}

\begin{corollary}[sharpness within the class]\label{cor:sharp}
For every \(n\ge5\), subdivide one edge of \(K_{1,3}\) exactly \(n-4\) times.
The resulting tree \(B_n\) satisfies
\[
 \mdim(B_n)=3=|V(B_n)|-\diam(B_n)+1.
\]
\end{corollary}

\begin{proof}
The arm lengths are \(1,1,n-3\), so Theorem~\ref{thm:intro} gives an
upper bound of 3.  A nonpath graph cannot have multiset dimension 1, and no
graph has multiset dimension 2~\cite{Simanjuntak2017}; hence the value is 3.
The diameter is \(n-2\).
\end{proof}

\section{Independent exact verification through order 15}

We used two exact implementations with different tree generators, distance
algorithms, and representation evaluators.

The discovery implementation uses NetworkX 3.4.2 to generate nonisomorphic
trees, computes integer distances by breadth-first search, and examines every
vertex subset in increasing cardinality.  Every reported witness is checked
both by literal sorting of its distances and by counting selected vertices in
exact-distance spheres.

The independent implementation uses no graph library and imports no discovery
code.  It recursively generates rooted unlabeled trees: a rooted tree is
represented by the sorted multiset of the codes of the rooted components
below its root.  After forgetting the root, free trees are deduplicated by the
standard center-based parenthesis encoding used in tree isomorphism
\cite{AhoHopcroftUllman1974}.  Distances are recomputed by Floyd--Warshall,
and every subset is checked through literal sorted distance multisets.  Thus
neither a timeout nor failure to find a witness is interpreted as a proof.

For the one-vertex tree we use the literal convention
\(\mdim(K_1)=0\), witnessed by the empty set.  The two programs agree on the
number of free trees and the complete multiset dimension distribution at every
order.  The result is shown in Table~\ref{tab:enumeration}; a dash means zero.

\begin{table}[ht]
\centering
\caption{Exact multiset dimensions of all trees through order 15.}
\label{tab:enumeration}
\small
\setlength{\tabcolsep}{3.7pt}
\begin{tabular}{r r r r r r r r r r r}
\toprule
\(n\)&trees&\(\infty\)&0&1&3&4&5&6&7&9\\
\midrule
1&1&--&1&--&--&--&--&--&--&--\\
2&1&--&--&1&--&--&--&--&--&--\\
3&1&--&--&1&--&--&--&--&--&--\\
4&2&1&--&1&--&--&--&--&--&--\\
5&3&1&--&1&1&--&--&--&--&--\\
6&6&2&--&1&3&--&--&--&--&--\\
7&11&4&--&1&5&1&--&--&--&--\\
8&23&9&--&1&11&2&--&--&--&--\\
9&47&20&--&1&23&3&--&--&--&--\\
10&106&48&--&1&53&2&2&--&--&--\\
11&235&109&--&1&110&6&9&--&--&--\\
12&551&261&--&1&235&22&31&1&--&--\\
13&1301&636&--&1&500&74&82&7&1&--\\
14&3159&1595&--&1&1093&224&211&27&7&1\\
15&7741&4038&--&1&2389&662&523&99&27&2\\
\bottomrule
\end{tabular}
\end{table}

In total, both implementations check all 13,188 nonisomorphic trees of order
at most 15 and find no counterexample to (C).  This is an exact finite result,
but it is not evidence sufficient to settle the unbounded conjecture.

As a separate check on the construction in
Lemma~\ref{lem:binary-resolves}, an implementation sharing neither finite-tree
enumerator tested all 90,311 capacity-valid ordered arm-length tuples with
3 through 12 arms and lengths at most 8.  It reconstructed every graph,
checked literal representations, and verified the diameter budget.  The
symbolic proof, rather than this bounded check, is the certificate for the
infinite family.

\section{Discussion}

Theorem~\ref{thm:intro} closes the gap between the necessary short-arm
capacity condition of Kono and Zhang~\cite{KonoZhang2021} and existence of an
ID-coloring for starlike trees.  It also verifies the diameter conjecture on a
class whose degree is unbounded and in which known constructions have
unbounded ID-number~\cite{KonoZhang2021}.  The key feature is that all branches
meet at one center: the distance-sum congruence recovers depth before the
shell polynomial recovers the branch word.

This separation no longer holds verbatim at a second branching vertex.
Indeed, simple landmark rules based only on a diametral path fail even for a
four-arm starlike tree, although the binary construction repairs that case.
For general trees, repeated rooted branches at several vertices interact, and
no global construction of size \(n-\diam(T)+1\) is presently obtained here.
The full conjecture therefore remains open.

\section*{Data, code, and computational disclosure}

The accompanying supplementary archive contains the LaTeX source, the two
independent finite enumerators, the fail-closed comparison program, the
independent starlike checker, and the JSON reports used for
Table~\ref{tab:enumeration}.  The computations used Python 3.10.12; only the
discovery implementation requires NetworkX 3.4.2.  All distances and subset
tests use exact integers.

Generative-AI tools assisted with code prototyping and language editing.  The
author independently checked the mathematical arguments, computations,
citations, and final manuscript and takes responsibility for them.

\bibliographystyle{amsplain}
\bibliography{references}

\end{document}
